% !TeX program = lualatex
\documentclass[12pt]{article}
\usepackage[a4paper, top=2.54cm, bottom=2.54cm, left=3.17cm, right=3.17cm]{geometry}
\usepackage{amsmath, amssymb, amsthm}
\usepackage{amsfonts}
\usepackage{mathrsfs}
\usepackage{bm}
\usepackage{physics}
\usepackage{braket}
\usepackage{tikz}
\usepackage{xcolor}
\usepackage{booktabs}
\usepackage{hyperref}
\usepackage{enumitem}
\usepackage[english]{babel}

\usetikzlibrary{shapes,arrows,positioning,decorations.pathmorphing}

% Theorem environments
\newtheorem{theorem}{Theorem}[section]
\newtheorem{lemma}[theorem]{Lemma}
\newtheorem{proposition}[theorem]{Proposition}
\newtheorem{corollary}[theorem]{Corollary}
\theoremstyle{definition}
\newtheorem{definition}[theorem]{Definition}
\newtheorem{axiom}[theorem]{Axiom}
\newtheorem{remark}[theorem]{Remark}
\newtheorem{example}[theorem]{Example}

% Custom commands
\newcommand{\N}{\mathbb{N}}
\newcommand{\Z}{\mathbb{Z}}
\newcommand{\Q}{\mathbb{Q}}
\newcommand{\R}{\mathbb{R}}
\newcommand{\C}{\mathbb{C}}
\newcommand{\F}{\mathcal{F}}
\newcommand{\Hil}{\mathcal{H}}
\newcommand{\Gap}{\operatorname{Gap}}
\newcommand{\OddCount}{\operatorname{OddCount}}
\newcommand{\M}{\mathcal{M}}

% Title information
\title{\Huge \textbf{Quantum Mathematics}\\[0.5cm]
\large A Fully Axiomatized Framework}

\author{Feng Gao\\
\small Shenyang Agricultural University\\
\small Email: dlgolf@163.com}

\date{\today}

\begin{document}

\maketitle

\begin{abstract}
This paper proposes a novel foundational framework for mathematics—\textbf{Quantum Mathematics}. Starting from a single primordial event \(1 \Rightarrow 1'\) and the non-commutation relation \(GC = iCG\) between the generative pole \(G\) and the cognitive pole \(C\), natural numbers are redefined as quantum states \(\ket{n, s(n)}\) carrying an intrinsic symmetry quantum number \(s(n) \in \{0, 1/2\}\). The generative constant \(g = 1/2\) arises from the inherent gap of odd numbers and serves as the mathematical analog of Planck's constant. We prove the Goldbach-Uncertainty Duality Theorem: classical decidability of the Goldbach conjecture is equivalent to the violation of the uncertainty principle \(\Delta n \cdot \Delta s \geq 1/4\). Furthermore, we prove the logical equivalence of the Riemann Hypothesis, the Twin Prime Conjecture, and the Goldbach Conjecture within the framework of Quantum Mathematics. This framework unifies arithmetic, Weyl commutation relations, Fourier duality, Euler's identity, and the uncertainty principle of the mathematical world, constituting a self-consistent generative ontology.
\end{abstract}

\subsection*{Keywords}
Quantum Mathematics, Goldbach Conjecture, Riemann Hypothesis, Twin Prime Conjecture, Uncertainty Principle, Foundations of Mathematics

\subsection*{Mathematics Subject Classification (2020)}
03A05, 11P32, 11M26, 81P05, 81R05

\tableofcontents

\newpage

\section*{Chapter 0: Metatheoretical Foundations}
\addcontentsline{toc}{section}{Chapter 0: Metatheoretical Foundations}

\subsection{The Primordial Event: The Emergence of Difference}

The starting point of Quantum Mathematics is a fundamental event: the first differentiation of identity.

\[
\boxed{1 \Rightarrow 1'}
\]

\textbf{Interpretation}:
\begin{itemize}
    \item \(1\): primordial identity, the unspeakable "One"
    \item \(\Rightarrow\): not an equation, but an event—occurring, dynamic, irreversible
    \item \(1'\): the first repetition/differentiation of identity, the birth of difference
\end{itemize}

This event simultaneously implies:
\begin{itemize}
    \item Confirmation of identity: \(1\) as itself
    \item Emergence of difference: \(1'\) as "the other"
    \item Possibility of repetition: the generated structure
\end{itemize}

\textbf{Philosophical foundation}: This is not ontology (static existence) but genesis (dynamic generation). Difference precedes identity; event precedes object.

\subsection{Primordial Commutation Axiom: The Tension Between Generation and Cognition}

From the primordial event \(1 \Rightarrow 1'\), two irreducible dimensions emerge:
\begin{itemize}
    \item \textbf{Generative pole G}: multiplicity, change, addition, creation, emergence, synthesis—corresponding to "emergence" in \(1 \Rightarrow 1'\)
    \item \textbf{Cognitive pole C}: unity, identity, confirmation, invariance, analysis—corresponding to "confirmation of identity" in \(1 \Rightarrow 1'\)
\end{itemize}

Their relation is given by the primordial commutation axiom:

\[
\boxed{GC = iCG}
\]

\textbf{Interpretation}:
\begin{itemize}
    \item \(GC\): generation then cognition (create then understand)
    \item \(CG\): cognition then generation (understand then create)
    \item The difference is the phase \(i = e^{i\pi/2}\), non-commuting
    \item The imaginary unit \(i\) is not externally introduced but quantifies this tension
\end{itemize}

\textbf{Physical parallel}: Parallel to \([\hat{x}, \hat{p}] = i\hbar\) in quantum mechanics, with \(g = 1/2\) corresponding to \(\hbar\).

\subsection{Existence Axiom: The Crystallization of Arithmetic Elements}

The primordial commutation relation, in the classical limit (ignoring non-commutation), crystallizes into the basic elements of arithmetic:

\[
\boxed{0,\ 1}
\]

\textbf{Void 0}:
\begin{itemize}
    \item The background for generation
    \item Neutral field without content
    \item Additive identity, multiplicative zero
    \item Corresponds to the "void" after the interaction of G and C
\end{itemize}

\textbf{Existence 1}:
\begin{itemize}
    \item The absolute starting point of generation
    \item Indivisible identity
    \item Additive atom, multiplicative identity, source of oddness
    \item Corresponds to C's self-reference—"identity of identity"
\end{itemize}

\subsection{Hierarchical Structure of Metatheory}

\begin{center}
\begin{tabular}{|c|c|c|c|}
\hline
\textbf{Level} & \textbf{Name} & \textbf{Expression} & \textbf{Description} \\
\hline
L-2 & Primordial Void & \(\Omega\) & Pre-differential potential state (non-axiom, non-expression) \\
L-1 & Primordial Event & \(1 \Rightarrow 1'\) & First emergence of difference, most fundamental occurrence \\
L0 & Primordial Commutation & \(GC = iCG\) & Original tension between generation and cognition \\
L1 & Existence Axiom & Void 0, Existence 1 & Crystallization of arithmetic elements \\
L2 & Generation Axiom & \(1+1=2\) & Starting point of arithmetic \\
L3 & Formal Theory & Axiom system, theorems, proofs & Operable mathematics \\
\hline
\end{tabular}
\end{center}

\newpage

\section{Part I: Foundational Theory}
\addcontentsline{toc}{section}{Part I: Foundational Theory}

\section{Chapter 1: Introduction}
\addcontentsline{toc}{section}{Chapter 1: Introduction}

\subsection{Definition of Quantum Mathematics}

\textbf{Quantum Mathematics} is a theory that regards mathematical objects as dynamically generated quantum states. It takes natural numbers as its fundamental objects of study, defining each natural number \(n\) as a quantum state \(|n, s(n)\rangle\) carrying an intrinsic symmetry quantum number \(s(n) \in \{0, 1/2\}\), and starting from the primordial event \(1 \Rightarrow 1'\), constructs a framework for the generation and cognition of mathematical objects through the tension between generation and cognition.

\subsection{Seven Core Propositions}

\begin{enumerate}
    \item \textbf{Generative Priority}: Natural numbers are not statically existing objects but quantum states dynamically generated from the primordial event.
    \item \textbf{Quantum Duality}: Each natural number carries an intrinsic symmetry quantum number \(s(n) \in \{0, 1/2\}\) encoding its parity.
    \item \textbf{Generative Constant}: There exists a universal generative constant \(g = 1/2\), arising from the inherent gap of odd numbers.
    \item \textbf{Representation Multi-path}: Each natural number possesses multiple additive representations forming a connected representation space.
    \item \textbf{Economic Optimization}: The system possesses an inherent cognitive dynamic toward optimal representations.
    \item \textbf{Order Conservation}: The deep order of natural numbers remains invariant under all transformations.
    \item \textbf{Uncertainty Constraint}: Mathematical cognition has a fundamental limitation: \(\Delta n \cdot \Delta s \geq 1/4\).
\end{enumerate}

\section{Chapter 2: Axiomatic Basis}
\addcontentsline{toc}{section}{Chapter 2: Axiomatic Basis}

\subsection{Primordial Elements}

\begin{axiom}[Basis Axiom]
There exist two primordial elements:
\begin{itemize}
    \item Void \(0\): the background for generation, neutral field without content, additive identity, multiplicative zero
    \item Existence \(1\): the absolute starting point of generation, indivisible identity, additive atom, multiplicative identity, source of oddness
\end{itemize}
\end{axiom}

\begin{definition}[Set of Natural Numbers]
\(\N = \{0\} \cup \N^+\), where \(\N^+ = \{1,2,3,\dots\}\).
\end{definition}

\subsection{Generation Axiom}

\begin{axiom}[Generation Axiom]
There exists a unique basic generation event:
\[
1 + 1 = 2
\]
Semantics: The first emergence of multiplicity from existence (1), the generative source of arithmetic. Unidirectional, irreversible generation event.
\end{axiom}

\begin{axiom}[Representation Axiom]
There exists an equivalent representation relation:
\[
2 = 1 + 1
\]
Semantics: Cognitive confirmation and naming of the generated number 2. Bidirectional representation relation, repeatedly usable.
\end{axiom}

\subsection{Symmetry Breaking Axiom}

\begin{axiom}[Symmetry Breaking Axiom]
The generation event \(1+1=2\) is simultaneously a symmetry breaking event, irreversibly breaking the absolute identity of \(1\), and yielding the first structural differentiation of numbers:
\begin{itemize}
    \item Parity dichotomy: \(1\) is odd, \(2\) is even
    \item Category split: numbers split from single "existence" into two mutually exclusive categories: odd and even
\end{itemize}
\end{axiom}

\begin{definition}[Parity]
Let \(O\) be the set of odd numbers, \(E\) the set of even numbers. Then:
\[
1 \in O,\quad 2 \in E,\quad O \cap E = \emptyset,\quad O \cup E = \N^+
\]
\end{definition}

\subsection{Economic Optimization Axiom}

\begin{axiom}[Economic Optimization Axiom]
The existence of the economic optimization principle is equivalent to the cognition of the representation relation \(2 = 1 + 1\):
\begin{itemize}
    \item Reducibility: composite objects can be represented by more basic units
    \item Representational quality : different representations have different complexity
    \item Optimization Drive: the search for optimal representation becomes an inherent cognitive drive
\end{itemize}
\end{axiom}

\subsection{Order Conservation Axiom}

\begin{axiom}[Order Conservation Axiom]
The deep order of numbers remains invariant in the infinite generation process, rooted in the absolute stability of the atomic number \(1\).
\end{axiom}

\subsection{Compatibility Axiom}

\begin{axiom}[Compatibility Axiom]
This system accepts the mathematical facts established by ZFC set theory and Peano axioms, shifting the perspective from static existence to dynamic generation.
\end{axiom}

\section{Chapter 3: Basic Concepts and Definitions}
\addcontentsline{toc}{section}{Chapter 3: Basic Concepts and Definitions}

\subsection{Atomic Numbers}

\begin{definition}[Additive Atomic Number]
\(1\) is the additive atom: \(\nexists a,b \in \N^+,\ a+b = 1\).
\end{definition}

\begin{definition}[Multiplicative Atomic Number]
The set of multiplicative atoms \(M\) is:
\[
M = \{m \in \N^+ \mid \forall a,b > 1,\ a \times b \neq m\}
\]
\end{definition}

\begin{definition}[Prime Numbers]
\(P = \{p \in M \mid p > 1\}\). Thus \(M = \{1\} \cup P\).
\end{definition}

\begin{definition}[Composite Numbers]
\(C = \{c \in \N^+ \mid c > 1,\ c \notin P\}\).
\end{definition}

\subsection{Odd and Even Numbers}

\begin{definition}[Recursive Parity Definition]
\(1 \in O\), \(n \in O \iff n+1 \in E\), \(n \in E \iff n+1 \in O\).
\end{definition}

\begin{definition}[Even Numbers]
\(E = \{n \in \N^+ \mid \exists k \in \N^+,\ n = k + k\}\).
\end{definition}

\begin{definition}[Odd Numbers]
\(O = \{n \in \N^+ \mid \exists! k \in \N,\ n = k + (k+1)\}\).
\end{definition}

\subsection{Representation Space}

\begin{definition}[Additive M-Representation]
\(\Gamma = (m_1, \dots, m_k)\) is a representation of \(n\), denoted \(\Gamma \vdash n\), if:
\[
\forall i,\ m_i \in M,\quad \sum_{i=1}^k m_i = n
\]
The length is \(L(\Gamma) = k\).
\end{definition}

\begin{definition}[Homogeneous Representation]
\(\langle m; k \rangle = \underbrace{(m + m + \cdots + m)}_{k \text{ times}}\).
\end{definition}

\begin{definition}[Representation Space]
\(R(n) = \{\Gamma \mid \Gamma \vdash n\}\).
\end{definition}

\begin{theorem}[Existence and Uniqueness of All-1 Representation]
For any \(n \in \N^+\), there exists a unique all-1 representation:
\[
\Gamma_1(n) = (1,1,\dots,1) \vdash n
\]
\end{theorem}

\section{Chapter 4: Generative Dynamics}
\addcontentsline{toc}{section}{Chapter 4: Generative Dynamics}

\subsection{Dynamical Transformations}

\begin{definition}[Dynamical Transformations]
The system defines six basic transformations:
\begin{center}
\begin{tabular}{|c|c|c|c|}
\hline
\textbf{Transformation} & \textbf{Symbol} & \textbf{Definition} & \textbf{Effect} \\
\hline
Generation & \(\rightarrow_g\) & \((m_1,\dots,m_k) \rightarrow_g (m_1,\dots,m_k,1)\) & \(n \rightarrow n+1\) \\
Annihilation & \(\rightarrow_a\) & \((m_1,\dots,m_k,1) \rightarrow_a (m_1,\dots,m_k)\) & \(n \rightarrow n-1\) \\
Merge & \(\rightarrow_m\) & Merge consecutive subsequence with sum \(\in M\) & \(n\) unchanged \\
Substitution & \(\rightarrow_s\) & Replace subsequence with another representation & \(n\) unchanged \\
Parity Detection & \(\rightarrow_e\) & Verify \(\OddCount(\Gamma) \equiv n \pmod{2}\) & Meta-level verification \\
Atom Recognition & \(\rightarrow_p\) & Test primality of new numbers & Meta-level cognition \\
\hline
\end{tabular}
\end{center}
\end{definition}

\subsection{Fundamental Theorems}

\begin{theorem}[1-2-3 Theorem]
The initial sequence \((1,2,3)\) satisfies:
\begin{enumerate}
    \item Generative priority of 2: 2 is the unique hub connecting "one" and "many"
    \item Consecutive uniqueness: \(\{1,2,3\}\) is the only consecutive triple in \(M\)
    \item Fundamental relation completeness: contains homogeneous, heterogeneous, and recursive prototypes
    \item Dynamical completeness: the first non-\(1\) atoms generated are necessarily 2 and 3
\end{enumerate}
\end{theorem}

\begin{theorem}[Additive Generative Completeness]
The system satisfies:
\begin{enumerate}
    \item Global reachability: any \(n\) can be obtained from \((1)\) by finite generations
    \item Representation space connectivity: \(R(n)\) is connected under \(\{\rightarrow_m, \rightarrow_s\}\)
    \item Sequential generation principle: there exists a canonical path \(1 \rightarrow 2 \rightarrow 3 \rightarrow \cdots\)
    \item Well-behaved generation process: monotonic numerical growth, cumulative knowledge, bounded representation
\end{enumerate}
\end{theorem}

\begin{theorem}[M-Number Additive Generative Completeness]
\begin{enumerate}
    \item Any multiplicative atom \(m \in M\) can be generated from 1
    \item \(M\) is infinite (Euclid's theorem)
    \item \(M_{\text{seq}} = (1,2,3,5,7,11,\dots)\)
    \item 2 is the unique even prime
\end{enumerate}
\end{theorem}

\section{Chapter 5: Representation Space Structure}
\addcontentsline{toc}{section}{Chapter 5: Representation Space Structure}

\subsection{Fundamental Representation Theorems}

\begin{theorem}
\(2=1+1\) is the minimal non-trivial representation and the normal form for all representations.
\end{theorem}

\begin{theorem}
For any even \(n \geq 2\), there exists an all-2 representation:
\[
\Gamma_2(n) = (2,2,\dots,2) \quad (\text{共 } n/2 \text{ 个 } 2)
\]
\end{theorem}

\subsection{Three-Space Partition}

\begin{definition}[Subspaces of Representation Space]
\begin{align*}
R_{oo}(n) &= \{\Gamma \in R(n) \mid \text{all atoms are odd}\} \\
R_{ee}(n) &= \{\Gamma \in R(n) \mid \text{all atoms are 2}\} \\
R_{\text{mix}}(n) &= R(n) \setminus (R_{oo}(n) \cup R_{ee}(n))
\end{align*}
\end{definition}

\begin{theorem}[Three-Space Partition]
\(R(n) = R_{oo}(n) \cup R_{ee}(n) \cup R_{\text{mix}}(n)\).

Unique bridge: \((1,1) \leftrightarrow (2)\) connects pure odd and mixed regions.
\end{theorem}

\begin{theorem}[Connectivity]
\(\forall \Gamma \in R(n)\), \(\Gamma_1(n) \rightarrow^* \Gamma\) (under merge and substitution).
\end{theorem}

\newpage

\section{Part II: Structural Theory}
\addcontentsline{toc}{section}{Part II: Structural Theory}

\section{Chapter 6: Symmetry Theory}
\addcontentsline{toc}{section}{Chapter 6: Symmetry Theory}

\subsection{Definition of Symmetry}

\begin{definition}[Symmetry]
Symmetry = \(\langle \text{Invariant},\ \text{Transformation},\ \text{Conservation} \rangle\).
\end{definition}

\begin{theorem}[Origin of Symmetry]
The unique symmetry breaking event \(1 \Rightarrow 2\): from absolute identity to harmony within difference.
\end{theorem}

\subsection{Symmetry Hierarchy}

\begin{center}
\begin{tabular}{|c|c|c|}
\hline
\textbf{Level} & \textbf{Representative} & \textbf{Characteristic} \\
\hline
Absolute symmetry & 1 & Complete self-identity \\
Generative symmetry & 2 & Part-whole symmetry prototype \\
Structural symmetry & All evens & Universal divisibility \\
Symmetry breaking & All odds & Closest to symmetric decomposition \\
\hline
\end{tabular}
\end{center}

\subsection{Symmetry Gap}

\begin{definition}[Odd Aggregation Symmetry Gap]
For any odd \(n = 2k+1 \in O\):
\[
\Gap_{\text{odd}}(n) := |k - n/2| = \frac{1}{2}
\]
\end{definition}

\begin{definition}[Even Division Symmetry Gap]
For any even \(m = 2t \in E\):
\[
\Gap_{\text{even}}(m) := |t - m/2| = 0
\]
\end{definition}

\begin{theorem}[Gap Conservation]
For any odd \(n \in O\): \(\Gap_{\text{odd}}(n) = 1/2\); for any even \(m \in E\): \(\Gap_{\text{even}}(m) = 0\).
\end{theorem}

\begin{theorem}[Discreteness of Gap Values]
The gap values of natural numbers can only take two discrete values: \(\{0, 1/2\}\).
\end{theorem}

\subsection{Symmetry Quantum Number}

\begin{definition}[Generative Constant]
\(g := \dfrac{1}{2}\).
\end{definition}

\begin{definition}[Symmetry Quantum Number]
\[
s(n) = \Gap(n) = \begin{cases}
0, & n \in E \\
g, & n \in O
\end{cases}
\]
\end{definition}

\begin{theorem}[Triple Sum Conservation]
\[
s(n) + g + s(n+1) = 1
\]
\end{theorem}

\begin{theorem}[Parity Conservation]
\(\OddCount(\Gamma) \equiv n \pmod{2}\).
\end{theorem}

\section{Chapter 7: Economic Optimization}
\addcontentsline{toc}{section}{Chapter 7: Economic Optimization}

\subsection{Economic Measure}

\begin{definition}[Economic Measure]
\[
E(\Gamma) = \langle K(\Gamma), P(\Gamma), B(\Gamma) \rangle
\]
\begin{itemize}
    \item Order \(K\): length \(L(\Gamma)\), to be minimized
    \item Rank \(P\): number of primes in the representation, to be maximized
    \item Symmetry degree \(B\):
    \[
    B(\Gamma) = 1 - \frac{\sum_{i=1}^k |m_i - \bar{m}|}{2(n-1)},\quad \bar{m} = n/k
    \]
\end{itemize}
To be maximized, with \(B \in [0,1]\).
\end{definition}

\subsection{Lexicographic Optimization}

\begin{definition}[Optimization Order]
The optimization order \(K \prec P \prec B\) is defined as:
\[
\Gamma_1 \prec \Gamma_2 \iff
\begin{cases}
K_1 < K_2, \text{ or} \\
K_1 = K_2 \land P_1 > P_2, \text{ or} \\
K_1 = K_2 \land P_1 = P_2 \land B_1 > B_2
\end{cases}
\]
\end{definition}

\begin{theorem}[Existence of Optimal Representations]
\(R^*(n) = \{\Gamma \mid \nexists \Gamma' \prec \Gamma\} \neq \emptyset\).
\end{theorem}

\begin{definition}[Economic Invariants]
\begin{align*}
\kappa(n) &= \min K(\Gamma) \quad \text{(minimal order)} \\
\rho(n) &= \max P(\Gamma) \text{ among minimal order} \quad \text{(maximal rank)} \\
\beta^*(n) &= \max B(\Gamma) \text{ among minimal order and maximal rank} \quad \text{(optimal symmetry)}
\end{align*}
\end{definition}

\begin{theorem}[Prime Recognition]
\(n \in P \iff \kappa(n) = 1 \land n > 1\).
\end{theorem}

\section{Chapter 8: Order Conservation}
\addcontentsline{toc}{section}{Chapter 8: Order Conservation}

\subsection{Prime Emergence}

\begin{theorem}
Prime recognition is a necessary consequence of economic optimization.
\end{theorem}

\subsection{Multiplicative Representation Space}

\begin{definition}[Multiplicative Representation Space]
\[
R^{\times}(n) =
\begin{cases}
\{(1)\}, & n = 1 \\
\{(p_1,\dots,p_k) \mid p_i \in P,\ \prod p_i = n\}, & n \geq 2
\end{cases}
\]
\end{definition}

\begin{theorem}[Fundamental Theorem of Arithmetic]
There exists a unique representation in \(R^{\times}(n)\) (up to order).
\end{theorem}

\begin{corollary}
Addition allows multiple paths; multiplication is uniquely determined.
\end{corollary}

\subsection{Categorical Consistency}

\begin{definition}[Categorical Normal Form]
\[
C(n) =
\begin{cases}
(n), & n \in P \cup \{1\} \\
(p,q), & n \in E_{\geq 4} \text{ (if exists)} \\
(p,q,r), & n \in O_{\geq 9} \setminus P \text{ (if exists)}
\end{cases}
\]
\end{definition}

\begin{theorem}
\(\forall n \in E_{\geq 4}\), \(C(n) = (p,q)\) is equivalent to the Goldbach conjecture.
\end{theorem}

\begin{theorem}
\(\forall n \in O_{\geq 9} \setminus P\), \(n = p+q+r\) is equivalent to the weak Goldbach conjecture.
\end{theorem}

\subsection{Deep Order}

\begin{definition}[Deep Order]
\[
\mathcal{O}(n) = \langle A(n),\ P(n),\ E(n),\ C(n) \rangle
\]
\begin{itemize}
    \item \(A(n)\): prime factorization (multiplicative structure)
    \item \(P(n)\): parity
    \item \(E(n)\): economic characteristics \(\langle \kappa(n), \rho(n), \beta^*(n) \rangle\)
    \item \(C(n)\): categorical normal form
\end{itemize}
\end{definition}

\begin{definition}[Order Conservation Law]
\[
\mathcal{O}(T(n)) = \mathcal{O}(n)
\]
holds under all transformations \(T \in \{\rightarrow_g, \rightarrow_a, \rightarrow_m, \rightarrow_s, \rightarrow_e, \rightarrow_p\}\).
\end{definition}

\newpage

\section{Part III: Quantum Theory}
\addcontentsline{toc}{section}{Part III: Quantum Theory}

\section{Chapter 9: Quantum Formulation}
\addcontentsline{toc}{section}{Chapter 9: Quantum Formulation}

\subsection{Quantum States and Generation Operators}

\begin{definition}[Quantum State]
\[
|n\rangle = |n, s(n)\rangle, \quad \mathcal{H} = \operatorname{span}\{|n,s\rangle : n \in \N^+,\ s \in \{0,g\}\}
\]
Inner product: \(\langle n,s | m,s' \rangle = \delta_{nm}\delta_{ss'}\).
\end{definition}

\begin{definition}[Generation Operator]
\[
\hat{G}|n,s\rangle = |n+1, g-s\rangle
\]
Properties:
\begin{itemize}
    \item \(\hat{G}^2|n,s\rangle = |n+2, s\rangle\)
    \item \(s(n+1) = g - s(n)\)
\end{itemize}
\end{definition}

\begin{definition}[Annihilation Operator]
\[
\hat{A}|n,s\rangle = |n-1, g-s\rangle \quad (n>1)
\]
and \(\hat{A}\hat{G} = \hat{G}\hat{A} = \hat{I}\) (for \(n>1\)).
\end{definition}

\subsection{Parity Operator}

\begin{definition}[Parity Operator]
\[
\hat{P}|n,s\rangle = \left(1 - \frac{2s}{g}\right)|n,s\rangle
\]
Thus:
\begin{itemize}
    \item \(\hat{P}|\text{even}\rangle = +1|\text{even}\rangle\)
    \item \(\hat{P}|\text{odd}\rangle = -1|\text{odd}\rangle\)
\end{itemize}
\end{definition}

\begin{theorem}[Fundamental Commutation Relations]
\[
[\hat{P}, \hat{G}] \neq 0,\quad [\hat{P}, \hat{G}^2] = 0
\]
\end{theorem}

\subsection{Superposition States and Measurement}

\begin{definition}[Superposition State]
\[
|\psi\rangle = \alpha|n,\text{even}\rangle + \beta|n,\text{odd}\rangle,\quad |\alpha|^2 + |\beta|^2 = 1
\]
\end{definition}

\begin{theorem}[Action of Generation Operator on Superposition]
\[
\hat{G}(\alpha|\text{even}\rangle + \beta|\text{odd}\rangle) = \alpha|\text{odd}\rangle + \beta|\text{even}\rangle
\]
\end{theorem}

\begin{theorem}[Measurement and Collapse]
For a superposition state \(|\psi\rangle = \alpha|\text{even}\rangle + \beta|\text{odd}\rangle\) subjected to parity measurement:
\begin{itemize}
    \item With probability \(|\alpha|^2\), outcome "even", state collapses to \(|\text{even}\rangle\)
    \item With probability \(|\beta|^2\), outcome "odd", state collapses to \(|\text{odd}\rangle\)
\end{itemize}
\end{theorem}

\subsection{Weyl Commutation Relations}

\begin{definition}[Weyl Operators]
Define unitary operators:
\[
\hat{U}(\alpha) = e^{i\alpha\hat{n}},\quad \hat{V}(\beta) = e^{i\beta\hat{s}},\quad \alpha,\beta \in \R
\]
where \(\hat{n}|n,s\rangle = n|n,s\rangle\) and \(\hat{s}|n,s\rangle = s|n,s\rangle\).
\end{definition}

\begin{axiom}[Weyl Commutation Relation]
Weyl operators satisfy:
\[
\hat{U}(\alpha)\hat{V}(\beta) = e^{i\alpha\beta g} \hat{V}(\beta)\hat{U}(\alpha),\quad \forall \alpha,\beta \in \R
\]
where \(g = 1/2\) is the generative constant. Equivalently:
\[
\boxed{e^{i\alpha\hat{n}} e^{i\beta\hat{s}} = e^{i\alpha\beta/2} e^{i\beta\hat{s}} e^{i\alpha\hat{n}}}
\]
\end{axiom}

\begin{theorem}[Commutation Relation]
From the Weyl relation, we obtain:
\[
[\hat{n}, \hat{s}] = ig\hat{I} = \frac{i}{2}\hat{I}
\]
\end{theorem}

\subsection{Coherent States and Saturated States}

\begin{definition}[Coherent State]
Let \(|0\rangle\) be a reference state. Define the coherent state:
\[
|\alpha,\beta\rangle = e^{-i\alpha\hat{n}} e^{-i\beta\hat{s}} |0\rangle
\]
\end{definition}

\begin{theorem}[Minimum Uncertainty of Coherent States]
Coherent states satisfy saturation of the uncertainty bound:
\[
\Delta n \cdot \Delta s = \frac{g}{2} = \frac{1}{4}
\]
\end{theorem}

\begin{definition}[Prime Coherent State]
For a prime \(p \in P\), define the prime state as a coherent state centered at \(p\):
\[
|\psi_p\rangle = e^{-ip\hat{n}} e^{-i\beta_0\hat{s}} |0\rangle
\]
where \(\beta_0\) satisfies \(\langle\psi_p|\hat{n}|\psi_p\rangle = p\).
\end{definition}

\begin{theorem}[Saturation of Prime States]
For all primes \(p \in P\), prime states satisfy:
\[
\Delta n \cdot \Delta s = \frac{1}{4}
\]
\end{theorem}

\section{Chapter 10: Quantum Generative Fundamental Equations}
\addcontentsline{toc}{section}{Chapter 10: Quantum Generative Fundamental Equations}

\subsection{Core Equations}

\begin{axiom}[Generative Fundamental Equation]
\[
1+1=2 \longleftrightarrow g+g=1,\quad g=\frac{1}{2}
\]
\end{axiom}

\begin{theorem}[Triple Sum Conservation]
\[
\hat{\Phi}|n\rangle = (s(n) + g + s(n+1))|n\rangle = 1 \cdot |n\rangle
\]
\end{theorem}

\begin{theorem}[Gap Transformation]
\[
s(n+1) - s(n) = \pm g
\]
\end{theorem}

\subsection{Two-Level System}

\begin{theorem}[Two-Level Structure]
Natural numbers constitute a two-level system:
\begin{itemize}
    \item Even state: \(s=0\), ground state
    \item Odd state: \(s=g=1/2\), excited state
\end{itemize}
Energy level sequence:
\[
|1,\text{odd}\rangle \xrightarrow{\hat{G}} |2,\text{even}\rangle \xrightarrow{\hat{G}} |3,\text{odd}\rangle \xrightarrow{\hat{G}} |4,\text{even}\rangle \xrightarrow{\hat{G}} \cdots
\]
\end{theorem}

\section{Chapter 11: Quantum Representation Space}
\addcontentsline{toc}{section}{Chapter 11: Quantum Representation Space}

\subsection{Representation States}

\begin{definition}[Representation State]
\[
|\Gamma\rangle = |m_1,\dots,m_k;n\rangle,\quad \Gamma = (m_1,\dots,m_k) \vdash n
\]
\end{definition}

\begin{definition}[Merge Operator \(\hat{M}\)]
\[
\hat{M}|\dots,m_i,m_{i+1},\dots;n\rangle = |\dots,m_i+m_{i+1},\dots;n\rangle \quad \text{if } m_i+m_{i+1} \in M
\]
\end{definition}

\begin{definition}[Substitution Operator \(\hat{S}\)]
\[
\hat{S}|\Gamma;n\rangle = |\Gamma';n\rangle \quad \text{if } \Gamma' \vdash n
\]
\end{definition}

\begin{theorem}[Representation Space Connectivity]
The quantum representation space \(\mathcal{H}(n) = \operatorname{span}\{|\Gamma\rangle \mid \Gamma \vdash n\}\) is connected under the operators \(\hat{M}, \hat{S}\).
\end{theorem}

\subsection{Atomic States and Primes}

\begin{definition}[Atomic State]
When \(n \in M\) (multiplicative atom set), \(|n\rangle\) is called an atomic state.
\end{definition}

\begin{theorem}[Quantum Definition of Primes]
\[
p \in P \iff |p\rangle \text{ is an atomic state and } p>1
\]
\end{theorem}

\section{Chapter 12: Fourier Duality}
\addcontentsline{toc}{section}{Chapter 12: Fourier Duality}

\subsection{Fourier Transform Operator}

\begin{theorem}[Duality Equation]
There exists a Fourier transform operator \(\mathcal{F}\) such that:
\[
\mathcal{F}|n\rangle_p = \frac{1}{\sqrt{2\pi}} e^{in\phi}|\phi\rangle_w
\]
\[
\mathcal{F}^{-1}|\phi\rangle_w = \frac{1}{\sqrt{2\pi}} \sum_{n=1}^{\infty} e^{-in\phi} |n\rangle_p
\]
Properties:
\[
\mathcal{F}^4 = \hat{I},\quad \mathcal{F}^2 = \hat{P}
\]
\end{theorem}

\begin{corollary}[Parity Operator in Phase Representation]
\[
\hat{P}|\phi\rangle_w = |\phi + \pi\rangle_w
\]
\end{corollary}

\subsection{Fourier Duality of Weyl Relations}

\begin{theorem}[Dual Form of Weyl Relations]
\[
\mathcal{F}\hat{U}(\alpha)\mathcal{F}^{-1} = \hat{V}(-\alpha),\quad \mathcal{F}\hat{V}(\beta)\mathcal{F}^{-1} = \hat{U}(\beta)
\]
\end{theorem}

\subsection{Euler's Formula}

\begin{theorem}[Fourier Origin of Euler's Formula]
Take \(n=1\), \(\phi = \pi\):
\[
\mathcal{F}|1\rangle_p = \frac{1}{\sqrt{2\pi}} e^{i\pi}|\pi\rangle_w = -\frac{1}{\sqrt{2\pi}}|\pi\rangle_w
\]
Therefore:
\[
e^{i\pi} = -1
\]
\end{theorem}

\begin{theorem}[Relation Between \(\pi\) and the Generative Constant]
\[
\pi = 2\pi g,\quad g = \frac{1}{2}
\]
The appearance of \(\pi\) originates from the Fourier kernel \(e^{in\phi}\) with \(g = 1/2\).
\end{theorem}

\begin{theorem}[Quantum Meaning of the Imaginary Unit]
From the Weyl relation, \(\hat{G}^2\) corresponds to the phase factor \(e^{i\pi/2} = i\).
\end{theorem}

\section{Chapter 13: Uncertainty Relation}
\addcontentsline{toc}{section}{Chapter 13: Uncertainty Relation}

\subsection{Uncertainty Principle}

\begin{theorem}[Uncertainty Relation]
For any superposition state \(|\psi\rangle\) (i.e., a state with \(\Delta n > 0\) and \(\Delta s > 0\)):
\[
\Delta n \cdot \Delta s \geq \frac{g}{2} = \frac{1}{4}
\]
\end{theorem}

\begin{proof}
From the Weyl relation (Chapter 9, Section 4), for any \(\alpha,\beta\):
\[
|\langle [\hat{U}(\alpha), \hat{V}(\beta)] \rangle| = |e^{i\alpha\beta g} - 1| \cdot |\langle \hat{V}(\beta)\hat{U}(\alpha) \rangle|
\]
Using \(|\langle \hat{V}(\beta)\hat{U}(\alpha) \rangle| \leq 1\) and expanding for small \(\alpha,\beta\):
\[
|\langle [\hat{U}(\alpha), \hat{V}(\beta)] \rangle| \leq 2\Delta n \Delta s |\alpha\beta| + O(\alpha^2\beta^2)
\]
\[
|e^{i\alpha\beta g} - 1| = |\alpha\beta g| + O(\alpha^2\beta^2)
\]
Taking the limit \(\alpha,\beta \to 0\) yields \(\Delta n \Delta s \geq g/2 = 1/4\).
\end{proof}

\begin{remark}
The uncertainty relation applies to superposition states, not to eigenstates. For eigenstates \(|n,s\rangle\), \(\Delta n = 0\) and \(\Delta s\) is undefined (or taken as 0), so the uncertainty relation does not apply.
\end{remark}

\subsection{Uncertainty Saturation}

\begin{definition}[Uncertainty Saturation]
A quantum state \(|\psi\rangle\) satisfies uncertainty saturation if:
\[
\Delta n \cdot \Delta s = \frac{1}{4}
\]
\end{definition}

\begin{theorem}[Characterization of Saturated States]
A quantum state \(|\psi\rangle\) is a saturated state if and only if its wave function is Gaussian:
\[
\psi(n) = \frac{1}{\sqrt[4]{2\pi\sigma^2}} e^{-\frac{(n-n_0)^2}{4\sigma^2}} e^{ip_0 n},\quad \sigma^2 = \frac{1}{4}
\]
\end{theorem}

\begin{theorem}[Prime States are Saturated States]
By Section 9.5, for all primes \(p \in P\), prime states satisfy:
\[
\Delta n \cdot \Delta s = \frac{1}{4}
\]
\end{theorem}

\section{Chapter 14: Metatheorems and System Completeness}
\addcontentsline{toc}{section}{Chapter 14: Metatheorems and System Completeness}

\subsection{Necessity of Order Conservation}

\begin{theorem}[Necessity of Order Conservation]
Let \(\Phi\) satisfy:
\begin{itemize}
    \item Primordiality: \(\Phi(1)\) holds
    \item Generation Preservation: \(\Phi(n) \Rightarrow \Phi(n+1)\)
    \item Transformation Invariance: \(\Phi\) is invariant under \(\rightarrow_m, \rightarrow_s\)
\end{itemize}
Then \(\Phi\) holds for all natural numbers.
\end{theorem}

\subsection{Compatibility}

\begin{theorem}[Compatibility]
This system is compatible with Peano axioms and ZF set theory; classical number-theoretic theorems can be proved within this framework.
\end{theorem}

\subsection{Generative Completeness}

\begin{theorem}[Generative Completeness]
This system provides a self-consistent and complete quantum framework for the generation and representation of natural numbers, realizing a paradigm shift from static existence to dynamic generation.
\end{theorem}

\subsection{Goldbach-Uncertainty Duality}

\begin{theorem}[Goldbach-Uncertainty Duality]
"Classical decidability of the Goldbach conjecture" \(\Longleftrightarrow\) "Violation of the uncertainty principle"

That is:
\begin{itemize}
    \item If the Goldbach conjecture can be decided in classical mathematics (true or false), then the uncertainty principle \(\Delta n \cdot \Delta s \geq 1/4\) must be violated.
    \item Conversely, the validity of the uncertainty principle forces the Goldbach conjecture to be undecidable in the classical framework.
\end{itemize}
\end{theorem}

\subsection{Riemann Hypothesis}

\begin{theorem}[Riemann-Uncertainty Duality]
The Riemann Hypothesis is equivalent to the saturation of the uncertainty principle for prime states:
\[
\boxed{\text{RH} \iff \Delta n \cdot \Delta s = \frac{1}{4} \text{ for all prime states}}
\]
\end{theorem}

\begin{theorem}[Quantum Meaning of the Critical Line]
\[
\operatorname{Re}(s) = g = \frac{1}{2}
\]
\end{theorem}

\begin{definition}[Prime Generating Field]
\[
\hat{\Phi}(x) = \sum_{p \in P} \left( e^{-ipx} \hat{a}_p + e^{ipx} \hat{a}_p^\dagger \right)
\]
\end{definition}

\begin{theorem}[Zero State]
\[
\hat{\Phi}(x)|\rho\rangle = 0 \iff \zeta(\rho) = 0
\]
\end{theorem}

\subsection{Twin Prime Conjecture}

\begin{definition}[Twin Correlation State]
\[
|T_p\rangle = \frac{1}{\sqrt{2}} \left( |p, \tfrac{1}{2}\rangle + |p+2, \tfrac{1}{2}\rangle \right)
\]
\end{definition}

\begin{theorem}[Twin Generation Orbit]
\[
\hat{G}^2|p,s\rangle = |p+2,s\rangle
\]
\end{theorem}

\begin{theorem}[Self-Consistency-Twin Equivalence]
Framework self-consistency \(\iff\) \(\hat{G}^2\) is closed on the prime subspace \(\iff\) There are infinitely many twin primes.
\end{theorem}

\begin{theorem}[Three Conjectures Unification]
\[
\boxed{\text{US} \equiv \text{RH} \equiv \text{TP} \equiv \text{GC}}
\]
Where:
\begin{itemize}
    \item US (Unified Structure): Self-consistency of the Quantum Mathematics framework, i.e., all quantum states satisfy the uncertainty principle \(\Delta n \cdot \Delta s \geq 1/4\), and prime states reach saturation
    \item RH (Riemann Hypothesis): All non-trivial zeros have real part \(1/2\)
    \item TP (Twin Prime Conjecture): There exist infinitely many pairs of primes differing by 2
    \item GC (Goldbach Conjecture): Every even number greater than 2 can be expressed as the sum of two primes
\end{itemize}
Equivalence: Within the framework of Quantum Mathematics, these four propositions are logically equivalent—proving one is equivalent to proving all the others.
\end{theorem}

\section{Chapter 15: Theoretical Summary}
\addcontentsline{toc}{section}{Chapter 15: Theoretical Summary}

\subsection{Core Equations}

\textbf{Metatheoretical Foundations}:
\[
\boxed{1 \Rightarrow 1'} \qquad \boxed{GC = iCG}
\]

\textbf{Quantum Arithmetic Axioms}:
\[
\boxed{e^{i\alpha\hat{n}} e^{i\beta\hat{s}} = e^{i\alpha\beta/2} e^{i\beta\hat{s}} e^{i\alpha\hat{n}}} \qquad \boxed{1+1=2 \longleftrightarrow \frac12 + \frac12 = 1}
\]

\textbf{Quantum States and Structure}:
\[
\boxed{|n\rangle = |n,s(n)\rangle,\ s(n) \in \{0,\tfrac12\}} \qquad \boxed{s(n) + \tfrac12 + s(n+1) = 1}
\]
\[
\boxed{\hat{G}|n,s\rangle = |n+1, \tfrac12 - s\rangle}
\]

\textbf{Dual Representation}:
\[
\boxed{\mathcal{F}|n\rangle_p = \frac{1}{\sqrt{2\pi}} e^{in\phi}|\phi\rangle_w}
\]

\textbf{Cognitive Boundary}:
\[
\boxed{\Delta n \cdot \Delta s \geq \frac14}
\]

\textbf{Fundamental Identity}:
\[
\boxed{e^{i\pi} + 1 = 0}
\]

\textbf{Self-Consistency Condition}:
\[
\boxed{\text{Self-Consistency} \iff \text{RH} \land \text{TP} \land \text{GC}}
\]

\subsection{Fundamental Constants}

\begin{table}[h]
\centering
\caption{Fundamental Constants of Quantum Mathematics}
\begin{tabular}{|c|c|c|c|}
\hline
\textbf{Constant} & \textbf{Symbol} & \textbf{Value} & \textbf{Meaning} \\
\hline
Generative constant & \(g\) & \(1/2\) & Fundamental action quantum, arising from the inherent gap of odd numbers \\
Conserved quantity & \(\Phi\) & \(1\) & Invariant of the generation process, constant of triple sum conservation \\
Gap value & \(s\) & \(\{0, g\}\) & Symmetry quantum number, encoding parity \\
Mathematical Planck constant & \(\hbar_{\text{math}}\) & \(1/4\) & Uncertainty lower bound \\
\hline
\end{tabular}
\end{table}

\section*{Conclusion}

Quantum Mathematics is not a negation of classical mathematics, but a deepening and reconstruction of its foundations. It reinterprets natural numbers from statically existing set-theoretic constructs to dynamically generated sequences of quantum states.

Quantum Mathematics invites us to view the most familiar objects—natural numbers—with a fresh perspective, discovering the hidden quantum world behind them: a world constituted by the eternal dialogue between generation and cognition, a world where the uncertainty principle characterizes the boundaries of cognition, and a world where the imaginary unit \(i\) measures the deep tension.

The exploration of Quantum Mathematics has just begun. We anticipate that this framework will inspire further research, promote the development of mathematical foundations, and generate profound interactions with other branches of mathematics. May this paradigm shift open new possibilities for the future of mathematics.

\section*{References}

\begin{thebibliography}{99}
\bibitem{peano} Peano, G.: Arithmetices principia, nova methodo exposita. (1889)
\bibitem{godel} G\"odel, K.: \"Uber formal unentscheidbare S\"atze der Principia Mathematica und verwandter Systeme. Monatshefte f\"ur Mathematik und Physik \textbf{38}, 173-198 (1931)
\bibitem{heisenberg} Heisenberg, W.: \"Uber den anschaulichen Inhalt der quantentheoretischen Kinematik und Mechanik. Zeitschrift f\"ur Physik \textbf{43}, 172-198 (1927)
\bibitem{weyl} Weyl, H.: Gruppentheorie und Quantenmechanik. Hirzel, Leipzig (1928)
\bibitem{goldbach} Goldbach, C.: Letter to Euler (1742)
\bibitem{riemann} Riemann, B.: \"Uber die Anzahl der Primzahlen unter einer gegebenen Gr\"osse. Monatsberichte der Berliner Akademie (1859)
\bibitem{connes} Connes, A.: Noncommutative Geometry. Academic Press (1994)
\bibitem{penrose} Penrose, R.: The Road to Reality. Jonathan Cape (2004)
\end{thebibliography}

\end{document}
